%
% connection_lease_reaping_destructive_read_buggy.tex
%
% FIDELITY CONTROL for H1 (destructive-read-sweep), not the design. Changes
% only Read: instead of Xi ConnReg (no change), it sweeps a lapsed
% connection out of `registered` -- exactly what HubClientRegistry
% ._sweep_locked does today on every call through named_sessions/
% live_sessions/repos -- and leaves `owner` untouched, because none of
% those six read call sites forwards the swept set to OwnerTracker. See
% docs/connection_lease_reaping.tex, Section "H1 --- destructive-read-sweep",
% for the full trace and the goal that must return FOUND against this
% variant and NOT found against the fixed spec.
%
% Type-check:  fuzz docs/connection_lease_reaping_destructive_read_buggy.tex
% Model-check:
%   probcli docs/connection_lease_reaping_destructive_read_buggy.tex \
%       -model_check -nodead -goal "..." -p DEFAULT_SETSIZE 3
%

\documentclass[a4paper,10pt,fleqn]{article}
\usepackage[margin=1in]{geometry}
\usepackage{fuzz}
\usepackage[colorlinks=true,linkcolor=blue,citecolor=blue,urlcolor=blue]{hyperref}

\begin{document}

\title{lux Connection-Lease Reaping --- Fidelity Control (Destructive Read)}
\author{Buggy variant: a read sweeps a lapsed connection without releasing it}
\date{September 2026}
\maketitle

% ============================================================
\section{Introduction}
% ============================================================

This is the fidelity control for H1 in
\texttt{docs/connection\_lease\_reaping.tex}, not the design. It removes
only \texttt{Read}'s non-destructiveness: a lapsed connection is still
correctly dropped from \texttt{registered} the moment some caller reads
through it, but every scene it owned is left exactly as it was --
matching \texttt{HubClientRegistry.named\_sessions}, which calls
\texttt{\_sweep\_locked} on every read and discards the swept set the
moment it returns, because none of its six callers
(\texttt{HubReads.client\_sessions}, \texttt{CallbackMenu.from\_named},
\texttt{client\_details.py}, \texttt{callback\_hold.py}, and the rest)
forwards that set to \texttt{OwnerTracker}. Read the main specification
first; this document states only what differs, in
Section~\ref{sec:read}.

% ============================================================
\section{Basic Types}
% ============================================================

\begin{zed}
  [CONN, IDENT, SCENE]
\end{zed}

Unchanged from the main specification.

% ============================================================
\section{Free Types}
% ============================================================

\begin{zed}
  TSTATE ::= tup | tdown
\end{zed}

\begin{zed}
  LSTATE ::= llive | llapsed
\end{zed}

\begin{zed}
  OWNER ::= unowned | conn \ldata CONN \rdata
\end{zed}

Unchanged from the main specification.

% ============================================================
\section{State}
% ============================================================

\begin{schema}{ConnReg}
  registered : \power CONN \\
  identity : CONN \pfun IDENT \\
  transport : CONN \pfun TSTATE \\
  lease : CONN \pfun LSTATE \\
  owner : SCENE \fun OWNER
\where
  registered \subseteq \dom identity \\
  \dom identity = \dom transport \\
  \dom transport = \dom lease
\end{schema}

Unchanged from the main specification.

% ============================================================
\section{Initialization}
% ============================================================

\begin{schema}{Init}
  ConnReg'
\where
  registered' = \emptyset \\
  identity' = \emptyset \\
  transport' = \emptyset \\
  lease' = \emptyset \\
  owner' = (\lambda s : SCENE @ unowned)
\end{schema}

Unchanged from the main specification.

% ============================================================
\section{Operations}
% ============================================================

\subsection{Read --- BUGGY: sweeps without releasing}
\label{sec:read}

\textbf{This is the one operation that differs from the design.} A
lapsed connection is still correctly dropped from \texttt{registered} --
the sweep itself is not the defect, and \texttt{Depart}/\texttt{Claim}'s
release mechanisms are deliberately left intact by this variant,
isolating the H1 finding from H2 and H3. But this happens from a query,
not the coordinator, and \texttt{owner} is left completely unchanged: no
scene the swept connection owned is released, and because the connection
is now outside \texttt{registered}, neither \texttt{Depart} nor
\texttt{Claim}'s sweep can ever find it again to release it later.

\begin{schema}{Read}
  \Delta ConnReg \\
  c? : CONN
\where
  c? \in registered \\
  lease~c? = llapsed \\
  registered' = registered \setminus \{ c? \} \\
  owner' = owner \\
  identity' = identity \\
  transport' = transport \\
  lease' = lease
\end{schema}

\subsection{Connect, Renew, TransportDies, LeaseLapse, Depart, Claim --- unchanged}

Identical to the main specification: only \texttt{Read} is no longer a
pure query.

\begin{schema}{Connect}
  \Delta ConnReg \\
  c? : CONN \\
  i? : IDENT
\where
  registered' = registered \cup \{ c? \} \\
  identity' = identity \oplus \{ c? \mapsto i? \} \\
  transport' = transport \oplus \{ c? \mapsto tup \} \\
  lease' = lease \oplus \{ c? \mapsto llive \} \\
  owner' = owner
\end{schema}

\begin{schema}{Renew}
  \Delta ConnReg \\
  c? : CONN
\where
  c? \in registered \\
  transport~c? = tup \\
  lease' = lease \oplus \{ c? \mapsto llive \} \\
  registered' = registered \\
  identity' = identity \\
  transport' = transport \\
  owner' = owner
\end{schema}

\begin{schema}{TransportDies}
  \Delta ConnReg \\
  c? : CONN
\where
  c? \in registered \\
  transport~c? = tup \\
  transport' = transport \oplus \{ c? \mapsto tdown \} \\
  registered' = registered \\
  identity' = identity \\
  lease' = lease \\
  owner' = owner
\end{schema}

\begin{schema}{LeaseLapse}
  \Delta ConnReg \\
  c? : CONN
\where
  c? \in registered \\
  lease~c? = llive \\
  lease' = lease \oplus \{ c? \mapsto llapsed \} \\
  registered' = registered \\
  identity' = identity \\
  transport' = transport \\
  owner' = owner
\end{schema}

\begin{schema}{Depart}
  \Delta ConnReg \\
  c? : CONN
\where
  c? \in registered \\
  registered' = registered \setminus \{ c? \} \\
  owner' = owner \oplus \\
  \quad~ (\{ s : SCENE | owner~s = conn(c?) \} \cross \{ unowned \}) \\
  identity' = identity \\
  transport' = transport \\
  lease' = lease
\end{schema}

\begin{schema}{Claim}
  \Delta ConnReg \\
  c? : CONN \\
  s? : SCENE
\where
  c? \in registered \\
  transport~c? = tup \\
  lease' = lease \oplus \{ c? \mapsto llive \} \\
  registered' = registered \setminus \\
  \quad~ \{ c2 : CONN | c2 \in registered \land c2 \neq c? \land
    lease~c2 = llapsed \} \\
  (owner \oplus (\{ s : SCENE | \exists c2 : CONN @ \\
    \qquad c2 \in registered \land c2 \neq c? \land lease~c2 = llapsed
    \land owner~s = conn(c2) \} \\
    \qquad \cross \{ unowned \}))~s? = unowned \\
  \lor \\
  (owner \oplus (\{ s : SCENE | \exists c2 : CONN @ \\
    \qquad c2 \in registered \land c2 \neq c? \land lease~c2 = llapsed
    \land owner~s = conn(c2) \} \\
    \qquad \cross \{ unowned \}))~s? = conn(c?) \\
  owner' = (owner \oplus (\{ s : SCENE | \exists c2 : CONN @ \\
    \qquad c2 \in registered \land c2 \neq c? \land lease~c2 = llapsed
    \land owner~s = conn(c2) \} \\
    \qquad \cross \{ unowned \})) \oplus \{ s? \mapsto conn(c?) \} \\
  identity' = identity \\
  transport' = transport
\end{schema}

% ============================================================
\section{Verification: the goal that must flip}
% ============================================================

Against this variant, I2/I5's negation must be \emph{found} --- the trace
in \texttt{connection\_lease\_reaping.tex}, Section H1, reaches it
directly: \texttt{Connect}(c1, i1); \texttt{Claim}(c1, s1);
\texttt{TransportDies}(c1); \texttt{LeaseLapse}(c1); \texttt{Read}(c1).

\begin{verbatim}
  # must be FOUND against this variant (NOT found against the fixed spec)
  probcli docs/connection_lease_reaping_destructive_read_buggy.tex \
    -model_check -nodead \
    -goal "#s.(#c.(s:SCENE & c:CONN & owner(s) = conn(c) &
                   not(c : registered)))" \
    -p DEFAULT_SETSIZE 3
\end{verbatim}

I1 must stay unaffected --- this variant leaves \texttt{Depart} and
\texttt{TransportDies} untouched, so the chain
\texttt{TransportDies};\texttt{Depart} still completes exactly as in the
fixed spec:

\begin{verbatim}
  # must remain FOUND -- this variant does not touch Depart or
  # TransportDies, only what a read does to a lapsed connection's
  # registration
  probcli docs/connection_lease_reaping_destructive_read_buggy.tex \
    -model_check -nodead \
    -goal "#c.(c:CONN & c : dom(identity) & not(c : registered) &
               transport(c) = tdown)" \
    -p DEFAULT_SETSIZE 3
\end{verbatim}

\end{document}
